\documentclass[11pt]{article}

% -------------------- Packages --------------------
\usepackage[margin=1in]{geometry}
\usepackage{amsmath, amssymb}
\usepackage{listings}
\usepackage{xcolor}
\usepackage{hyperref}
\usepackage{graphicx}
\usepackage{enumitem}
\usepackage{fancyhdr}
\usepackage[listings]{tcolorbox}
\tcbuselibrary{breakable}
\usepackage{colortbl}
\usepackage{tikz}
\usetikzlibrary{arrows.meta,calc}

% ---------- Modular row commands ----------
\newcommand{\memrow}[1]{\texttt{#1} \\ \hline}
\newcommand{\highlightmemrow}[1]{\cellcolor{gray!20}\texttt{#1} \\ \hline}

% ---------- Memory table environment ----------
\newenvironment{memorytable}[1][4cm]{
\begin{center}
\begin{tabular}{|p{#1}|}
\hline
}{
\end{tabular}
\end{center}
}

% ---------- Memory diagram (address -> value) ----------
\newcommand{\memcell}[2]{\texttt{#1} & \texttt{#2} \\ \hline}
\newcommand{\highlightmemcell}[2]{\cellcolor{gray!20}\texttt{#1} & \cellcolor{gray!20}\texttt{#2} \\ \hline}

% Named cells (for drawing arrows between rows/cells)
\newcommand{\memmark}[2]{\tikz[remember picture,baseline=(#1.base)]\node[inner sep=0pt,outer sep=0pt] (#1) {#2};}
\newcommand{\memcellm}[4]{\memmark{#1}{\texttt{#2}} & \memmark{#3}{\texttt{#4}} \\ \hline}
\newcommand{\highlightmemcellm}[4]{\cellcolor{gray!20}\memmark{#1}{\texttt{#2}} & \cellcolor{gray!20}\memmark{#3}{\texttt{#4}} \\ \hline}

% Draw an arrow from one named cell to another.
\newcommand{\memarrow}[3][bend left=20]{%
\tikz[remember picture,overlay] \draw[-{Stealth[length=2mm,width=2mm]}] (#2.east) to[#1] (#3.east);
}

\newenvironment{memorydiagram}[1][Memory Diagram]{
\begin{tcolorbox}[
    breakable,
    colback=gray!5,
    colframe=gray!50,
    title={#1}
]
\begin{center}
\begin{tabular}{|p{3.5cm}|p{9cm}|}
\hline
\textbf{Address} & \textbf{Value} \\ \hline
}{
\end{tabular}
\end{center}
\end{tcolorbox}
}

% -------------------- Header / Footer --------------------
\pagestyle{fancy}
\fancyhf{}
\lhead{CSI2110 Data Structures \hfill Nathan Ng}
\cfoot{\thepage}

% -------------------- Code Styling --------------------
\lstset{
    basicstyle=\ttfamily\small,
    keywordstyle=\color{blue},
    commentstyle=\color{gray},
    stringstyle=\color{teal},
    numbers=left,
    numberstyle=\tiny,
    stepnumber=1,
    numbersep=8pt,
    frame=single,
    breaklines=true,
    tabsize=2,
    language=Java
}

% -------------------- Custom Commands --------------------
\newcommand{\topic}[1]{\section*{#1}}
\newcommand{\subtopic}[1]{\subsection*{#1}}
\newcommand{\important}{\textbf{Important: }}

% ---------- Definitions (boxed + auto-numbered) ----------
\newcounter{definition}
\newenvironment{definition}[1][]
{%
    \refstepcounter{definition}%
    \begin{tcolorbox}[
        breakable,
        colback=gray!5,
        colframe=gray!50,
        title={Definition~\thedefinition\if\relax\detokenize{#1}\relax\else\ (#1)\fi}
    ]%
}
{%
    \end{tcolorbox}%
}

\newcommand{\defn}[2][]{\begin{definition}[#1]#2\end{definition}}

% ---------- Big Questions (boxed + auto-numbered) ----------
\newcounter{bigquestion}
\newenvironment{bigquestion}[1][]
{%
    \refstepcounter{bigquestion}%
    \begin{tcolorbox}[
        breakable,
        colback=gray!5,
        colframe=gray!50,
        title={Big Question~\thebigquestion\if\relax\detokenize{#1}\relax\else\ (#1)\fi}
    ]%
}
{%
    \end{tcolorbox}%
}

\newcommand{\bigq}[2][]{\begin{bigquestion}[#1]#2\end{bigquestion}}
\newcommand{\bigqa}[3][]{%
    \begin{bigquestion}[#1]
        \textbf{Q:} #2\par\medskip
        \textbf{A:} #3
    \end{bigquestion}%
}

% -------------------- Document --------------------
\begin{document}

\begin{center}
    {\Large \textbf{CSI2110 -- Review Notes}}\\
\end{center}

\topic{Arrays}
\defn[Array Memory Layout]{
In memory, array elements are stored contiguously. For an \texttt{int[]} in Java, each element is typically 4 bytes apart.
}

\begin{lstlisting}
int[] arr = {10, 20, 30, 40};
\end{lstlisting}

\begin{memorydiagram}[Array Example in Memory]
  \memcell{1000}{arr[0] = 10}
  \memcell{1004}{arr[1] = 20}
  \memcell{1008}{arr[2] = 30}
  \memcell{1012}{arr[3] = 40}
\end{memorydiagram}

\topic{Singly Linked Lists (SLList)}
\defn[Singly Linked List]{
Each node stores a value and a reference to the next node. Nodes are not required to be contiguous in memory.
}

\begin{memorydiagram}[Singly Linked List Node Layout]
  \memcell{1000}{[10 | next $\to$ 5000]}
  \memcell{5000}{[20 | next $\to$ 2400]}
  \memcell{2400}{[30 | next $\to$ null]}
\end{memorydiagram}

\subtopic{Operation: Insert at Position \texttt{i}}
Insert a new node so that:
\begin{itemize}
  \item The new node points to the old successor node.
  \item The previous node now points to the new node.
\end{itemize}

\begin{lstlisting}
void insertAt(int i, int value) {
    Node prev = head;
    for (int k = 0; k < i - 1; k++) {
        prev = prev.next;
    }
    Node n = new Node(value);
    n.next = prev.next;
    prev.next = n;
}
\end{lstlisting}

\bigqa[Tail Pointer in SLList]
{Suppose we want to quickly insert or delete the last node. Is it sufficient to store only a tail pointer?}
{A tail pointer helps with inserting at the end in $O(1)$ when we already know the tail. However, deleting the last node still requires locating the second-last node, so in a singly linked list deletion at the end is not $O(1)$ unless extra structure is maintained.}

\topic{Doubly Linked Lists (DLList)}
\defn[Doubly Linked List]{
A doubly linked list stores \texttt{previous} and \texttt{next} references in each node, allowing two-way traversal.
}

\begin{memorydiagram}[Doubly Linked List Node Layout]
  \memcell{1000}{[null $\leftarrow$ previous | 10 | next $\to$ 5000]}
  \memcell{5000}{[1000 $\leftarrow$ previous | 20 | next $\to$ 2400]}
  \memcell{2400}{[5000 $\leftarrow$ previous | 30 | next $\to$ null]}
\end{memorydiagram}

\subtopic{Operation: Delete Node at Position \texttt{i}}
Deleting node \texttt{i} means:
\begin{itemize}
  \item Find the target node.
  \item Redirect the previous node's \texttt{next} pointer.
  \item Redirect the next node's \texttt{previous} pointer.
\end{itemize}

\begin{lstlisting}
void deleteAt(Node x) {
    if (x.prev != null) {
        x.prev.next = x.next;
    }
    if (x.next != null) {
        x.next.prev = x.prev;
    }
}
\end{lstlisting}

\important{In DLList, delete is cleaner because each node already has direct access to both neighbors.}

\end{document}
